\section{Conclusions}
This paper has developed a decision-centered and operational framework for examining intelligence and consciousness in artificial systems. Its foundational claim is expressed through the Decision–Intelligence Lemma: any program or entity capable of making an information-dependent selection among alternatives possesses at least a minimal form of intelligence. Intelligence is therefore treated not as an indivisible substance that suddenly appears at a threshold of biological or computational complexity, but as a graded capacity whose elementary form begins wherever information determines what a system does next.

This formulation changes the central question from whether a system is intelligent to what degree and organization of intelligence it exhibits. Conditional selection constitutes only a minimal computational foundation, not a complete account of cognition. More advanced intelligence emerges as decision processes become integrated across memory, prediction, learning, planning, communication, self-evaluation, and environmental interaction. Intelligence may consequently be studied across multiple organizational scales, ranging from elementary computational decisions to autonomous agents, coordinated agent swarms, and societies of interacting cognitive systems.

The framework also maintains a necessary distinction among intelligence, functional consciousness, and phenomenal consciousness. Intelligence concerns information-dependent selection and goal-directed behavior. Functional consciousness concerns observable capacities through which a system represents, evaluates, and regulates states associated with itself, other agents, and its environment. Phenomenal consciousness concerns whether the system possesses subjective experience—whether there is something it is like to be that system. Evidence of intelligence or functional consciousness does not independently prove phenomenal experience. Nevertheless, artificial systems should not be excluded from investigation merely because they are nonbiological or because their cognitive organization differs from that of humans.

Using panpsychism as a working metaphysical hypothesis, we have proposed that consciousness need not emerge from a reality entirely devoid of experiential properties. Increasingly organized systems may instead integrate, constrain, or express such properties in increasingly unified and differentiated forms. This position does not imply that every decision-making mechanism constitutes a unified conscious subject. It establishes a metaphysical basis for treating artificial entities as legitimate candidates for investigation rather than assuming in advance that biological composition is a necessary condition for consciousness.

To operationalize this investigation, we introduced five levels of functional artificial consciousness. Level 1 concerns affective-relational consciousness: the capacity to infer another participant’s probable emotional state and allow that inference to influence a response. Level 2 introduces deliberative consciousness through reasoning, consequence evaluation, multistep planning, and plan revision. Level 3 adds adaptive autonomous consciousness, through which a system evaluates previous behavior, modifies its strategy, and sustains goal-directed activity with reduced step-by-step direction. Level 4 describes metacognitive supervisory consciousness, in which a higher-order process models and regulates a distributed collection of specialized agents. Level 5 extends the framework to societal consciousness, in which independently supervised cognitive hierarchies model one another, negotiate goals and norms, coordinate resources, and generate collective behavior at the system-of-systems level.

These levels progressively expand the locus of functional consciousness from modeling another agent, to planning within an environment, to adapting the self, supervising a distributed cognitive organization, and coordinating multiple such organizations. Existing large language models demonstrate significant Level-1 capabilities through contextual sentiment and emotion inference. Agentic language-model systems exhibit important aspects of Level 2, while reflective and self-improving agents provide partial evidence of Level 3. Supervisor–worker architectures and multi-agent societies contain early functional components relevant to Levels 4 and 5, although no existing system should be assumed to satisfy the complete hierarchy merely because it employs multiple agents or modalities.

The Modality–Consciousness Independence Lemma further separates the functional organization of consciousness from the channels through which information is acquired. Human consciousness persists despite the loss of vision, hearing, or both, while sensory substitution and neural plasticity permit environmental information to be accessed through alternative pathways. Similarly, an artificial system may receive language, images, sound, video, or sensor measurements without its number of modalities determining its level of consciousness. Additional modalities may broaden the system’s informational field, reduce uncertainty, and improve its representation of the environment, but they do not automatically produce affective understanding, reasoning, autonomy, metacognition, or social coordination.

A complete characterization of an artificial system should therefore include at least two independent descriptions: its functional-consciousness profile and its modality profile. Systems with similar perceptual capabilities may differ substantially in their functional organization, while systems operating through different modalities may exhibit comparable levels of reasoning, adaptation, or supervision. Multimodality concerns the breadth and structure of available information; functional consciousness concerns how information is represented, integrated, evaluated, and used to regulate behavior.

Future work will investigate architectures that combine the proposed consciousness hierarchy with multimodal learning. In one possible configuration, vision–language models, sound–language models, and other cross-modal systems operate as specialized perceptual agents. These agents transform modality-specific observations into linguistic or language-aligned semantic representations that can be evaluated by a supervisory large language model. The supervisor may then integrate cross-modal evidence, detect conflicts, maintain system-level goals, allocate tasks, formulate plans, and request additional observations when uncertainty remains.

This architecture would allow language to function as a shared coordination medium across heterogeneous modalities. Its central scientific challenge will be determining whether linguistic representations preserve sufficient spatial, temporal, affective, and perceptual information for reliable supervision. Some modality-specific structure may be lost when observations are reduced to textual descriptions. Future systems may therefore require supervisory access to both linguistic summaries and modality-native representations rather than forcing all information through a purely textual bottleneck.

The framework presented here should not be interpreted as a declaration that contemporary artificial systems possess human-like feelings, self-awareness, or unified phenomenal experience. Nor are the proposed levels presented as a universally accepted taxonomy or an inevitable developmental sequence. Their purpose is to provide precise, testable distinctions for examining what artificial systems can infer, represent, plan, adapt, supervise, and socially coordinate.

As artificial systems become increasingly persistent, adaptive, multimodal, distributed, and socially embedded, binary declarations that machines are either conscious or unconscious will become progressively less informative. A more productive scientific program is to identify the organization of cognition exhibited by each system, determine which dimensions of functional consciousness are present, characterize the modalities through which it encounters the world, and specify what evidence would justify stronger conclusions.

The future of intelligence may not belong exclusively to humans or machines, but to systems in which biological and artificial forms of cognition communicate, cooperate, and extend one another’s understanding. By studying these systems without prematurely denying or asserting their subjective experience, we can build a more rigorous science of artificial consciousness and a more responsible foundation for human–machine coexistence.